%% ch03.tex --- Gradient Descent
%% Numbers are computed by verify/ch03_gradient_descent.py.

\chapter{Gradient Descent}
\label{ch:ch03}

\epigraph{An algorithm must be seen to be believed.}{---\ Donald E.\ Knuth}

\begin{chapterlearning}
After working through this chapter you will be able to:
\begin{enumerate}
  \item State the gradient descent update rule and explain why the derivative
        gives both the direction and the size of the step.
  \item Perform the update \emph{simultaneously}, and say what different
        algorithm you get if you do not.
  \item Derive $\partial J/\partial w$ and $\partial J/\partial b$ for the
        squared-error cost, and implement the resulting algorithm.
  \item Diagnose a learning rate that is too small or too large from a plot of
        $J$ against iteration, and derive the exact threshold at which gradient
        descent on a quadratic cost begins to diverge.
  \item Carry out gradient descent by hand for two iterations and check your
        arithmetic against a printed table.
  \item Explain why the squared-error cost of a linear model is convex, and
        what that guarantees.
  \item Choose a stopping rule, and explain why a small change in $J$ is weak
        evidence of convergence.
  \item Distinguish batch, stochastic and mini-batch gradient descent, and say
        when each is appropriate.
\end{enumerate}
\end{chapterlearning}

\section{The Problem: Minimising Without Solving}
\label{sec:minimising}

\Cref{ch:ch02} ended by minimising $J$ exactly. Set both partial derivatives to
zero, solve the resulting pair of equations, and out comes the best line in one
step. Why, then, an entire chapter on an iterative method that only approaches
the answer?

Because the closed form is a happy accident of this particular model. It exists
because $J$ is a quadratic in the parameters, so its derivatives are linear and
the stationarity equations are a linear system. Change almost anything --- use
the logistic model of \cref{ch:ch06}, stack layers as in \cref{ch:ch09}, add a
non-quadratic penalty --- and setting the derivatives to zero produces equations
that nobody can solve symbolically. Even for linear regression, the closed form
requires inverting a matrix whose size is the number of features, which becomes
impractical well before the interesting problems begin (\cref{ch:ch04}).

So we need a method that works with much less. Assume the only two things we can
ever do are:

\begin{itemize}
  \item evaluate $J$ at a point $(w,b)$, and
  \item evaluate its derivatives at that point.
\end{itemize}

\noindent We cannot see the whole surface of \cref{fig:bowl}; we can only stand
at one point of it and feel which way is downhill. The algorithm of this chapter
is the obvious response to that situation, and it is, with refinements, how
essentially every model in this book is trained.

\begin{keyinsight}{The bargain}
Gradient descent trades exactness for generality. It never solves anything, and
it never arrives --- it only gets closer. In exchange it needs nothing from the
cost function except that we can differentiate it, which is why the same twenty
lines of code train a straight line and a neural network.
\end{keyinsight}

\section{Derivative Intuition in One Dimension}
\label{sec:derivative-intuition}

Fix $b$ and consider $J$ as a function of $w$ alone, so that the picture is the
parabola of \cref{sec:cost-one-param}. Take the three-point set of that section,
for which
\[
  J(w) = \tfrac{14}{6}(w-1)^2,
  \qquad
  \frac{\dd J}{\dd w} = \tfrac{14}{3}(w-1),
\]
with its minimum at $w = 1$. Stand at some $w$ and ask which way to move.

\begin{figure}[htbp]
  \centering
  \begin{tikzpicture}
    \begin{axis}[mlaxis, width=0.72\linewidth, height=5.4cm,
      xlabel={$w$}, ylabel={$J(w)$},
      xmin=-0.6, xmax=2.6, ymin=-0.6, ymax=5.4,
      legend style={font=\tiny, at={(0.5,0.98)}, anchor=north}]
      \addplot[MLBlue,very thick,domain=-0.6:2.6,samples=90]{14/6*(x-1)^2};
      \addlegendentry{$J(w)$}
      % tangent at w = 2 : slope +14/3
      \addplot[MLRose,thick,dashed,domain=1.35:2.55,forget plot]
        {2.3333 + 4.6667*(x-2)};
      % tangent at w = 0 : slope -14/3
      \addplot[MLTeal,thick,dashed,domain=-0.55:0.65,forget plot]
        {2.3333 - 4.6667*(x-0)};
      \addplot[only marks,mark=*,mark size=2pt,MLRose,forget plot]
        coordinates {(2,2.3333)};
      \addplot[only marks,mark=*,mark size=2pt,MLTeal,forget plot]
        coordinates {(0,2.3333)};
      \draw[-{Stealth},MLRose,thick]  (axis cs:2,1.75) -- (axis cs:1.533,1.75);
      \draw[-{Stealth},MLTeal,thick]  (axis cs:0,1.75) -- (axis cs:0.467,1.75);
      \node[font=\tiny,MLRose,anchor=south east] at (axis cs:2.5,3.0)
        {$J'>0$: move left};
      \node[font=\tiny,MLTeal,anchor=south west] at (axis cs:-0.55,3.0)
        {$J'<0$: move right};
      \node[font=\tiny,MLBlue,anchor=north] at (axis cs:1,-0.05) {$J'=0$};
    \end{axis}
  \end{tikzpicture}
  \caption{The derivative points \emph{away} from the minimum, so we move
  against it. At $w=2$ the slope is $+4.667$ and we step left; at $w=0$ it is
  $-4.667$ and we step right. Both steps have length $\alpha\cdot 4.667$, and
  both land closer to $w=1$.}
  \label{fig:tangents}
\end{figure}

At $w = 2$ the derivative is $+4.667$: the function is increasing to the right,
so decreasing to the left, so we should decrease $w$. At $w = 0$ it is $-4.667$
and we should increase $w$. In both cases we move in the direction
\emph{opposite} to the sign of the derivative, which is captured by subtracting
it:
\[
  w \;:=\; w - \alpha\,\frac{\dd J}{\dd w}.
\]
With $\alpha = 0.1$, the point at $w = 2$ moves to $2 - 0.1(4.667) = 1.533$ and
the point at $w = 0$ moves to $0 + 0.1(4.667) = 0.467$. Both are closer to the
minimum than they were.

Notice what else the derivative supplied. Its \emph{magnitude} is large where
the parabola is steep and small near the bottom, so the steps are large when far
away and shrink automatically as the minimum is approached. At $w = 1.6$ the
derivative is only $2.8$ and the step is correspondingly shorter. Nothing in the
algorithm has to detect proximity to the minimum: the gradient does it.

\begin{keyinsight}{What one derivative buys}
A single number, $\dd J/\dd w$, tells you which way to move (its sign) and
roughly how far (its magnitude). This is the entire content of gradient descent.
Everything else in this chapter is bookkeeping over several parameters, or
care about the constant $\alpha$.
\end{keyinsight}

\section{The Update Rule}
\label{sec:update-rule}

With two parameters, the same reasoning applies to each in turn, using partial
derivatives.

\begin{defbox}{Gradient descent}
Given a differentiable cost $J(w,b)$, a \term{learning rate} $\alpha>0$ and a
starting point, repeat until convergence:
\[
  \begin{aligned}
    w &:= w - \alpha\,\frac{\partial J}{\partial w}(w,b),\\[2pt]
    b &:= b - \alpha\,\frac{\partial J}{\partial b}(w,b),
  \end{aligned}
\]
where both right-hand sides are evaluated at the \emph{current} $(w,b)$ and the
two assignments take effect together. The symbol $:=$ denotes assignment, not
an assertion of equality.
\end{defbox}

That last requirement --- both gradients evaluated before either parameter
changes --- is called the \term{simultaneous update}, and it is the first thing
that goes wrong in a hand-written implementation. The correct order is: compute
both partial derivatives into temporary variables, then assign.

\begin{lstlisting}[style=pythonml,caption={Simultaneous update: compute both
gradients before assigning either parameter.}]
tmp_w = w - alpha * dJ_dw(w, b)     # both use the OLD w and b
tmp_b = b - alpha * dJ_db(w, b)
w, b = tmp_w, tmp_b                 # now assign
\end{lstlisting}

\begin{pitfall}{The non-simultaneous update is a different algorithm}
Update $w$ first, and then compute $b$'s gradient at the \emph{new} $w$, and you
are no longer doing gradient descent. You are doing something else --- a
coordinate-wise method related to Gauss--Seidel iteration.

What makes this dangerous is that it does not announce itself. On the four-point
set with $\alpha=0.1$ the two versions run like this:

\begin{center}
\small
\begin{tabular}{@{}rcccc@{}}
\toprule
& \multicolumn{2}{c}{simultaneous (correct)} & \multicolumn{2}{c}{sequential (wrong)}\\
\cmidrule(lr){2-3}\cmidrule(lr){4-5}
step & $(w,b)$ & $J$ & $(w,b)$ & $J$\\
\midrule
$1$ & $(1.0750,\ 0.3500)$ & $0.3042$ & $(1.0750,\ 0.0813)$ & $0.4646$\\
$2$ & $(1.2563,\ 0.3963)$ & $0.0995$ & $(1.3234,\ 0.0923)$ & $0.0777$\\
$3$ & $(1.2900,\ 0.3926)$ & $0.0920$ & $(1.3828,\ 0.0873)$ & $0.0555$\\
\bottomrule
\end{tabular}
\end{center}

The cost falls in both, and after three steps the buggy version has the
\emph{lower} cost. Watching $J$ decrease will therefore never reveal this
mistake. What you have lost is everything you can prove: the convergence
guarantee, the stability threshold of \cref{sec:learning-rate}, and any
correspondence between your code and the mathematics you derived. On a
problem where the two disagree --- and there are many --- you would have no way
to tell which of your reasoning and your program was at fault.
\end{pitfall}

\section{The Learning Rate}
\label{sec:learning-rate}

The learning rate $\alpha$ is the one number gradient descent cannot choose for
itself, and choosing it badly is the most common reason a first implementation
appears not to work.

\subsection{Too small}

If $\alpha$ is small, every step is short and progress is slow. It is not
\emph{wrong} --- the cost still decreases monotonically, and given enough
iterations it arrives --- it is merely expensive. With $\alpha = 0.01$ on the
four-point set, the cost is still $0.076$ after $120$ iterations, on its way to
a minimum of $0.025$.

\subsection{Too large}

If $\alpha$ is large, each step overshoots the minimum and lands on the far
side. If it overshoots by less than it started with, the process still
converges, in a zig-zag. If it overshoots by more, the distance grows every
step, $J$ increases without bound, and the parameters overflow to infinity in a
few dozen iterations.

\begin{figure}[htbp]
  \centering
  \begin{minipage}[t]{0.48\linewidth}
    \centering
    \begin{tikzpicture}
      \begin{axis}[mlaxis, width=\linewidth, height=4.8cm, ymode=log,
        xlabel={iteration}, ylabel={$J$},
        xmin=0, xmax=120, ymin=0.02, ymax=10,
        legend style={font=\tiny, at={(0.98,0.98)}, anchor=north east}]
        \addplot[MLTeal,thick] coordinates {
          (1,6.5252)(2,5.4964)(3,4.6322)(4,3.9063)(6,2.7843)(8,1.9925)
          (11,1.2195)(15,0.6543)(20,0.3278)(27,0.1615)(36,0.1050)
          (48,0.0901)(64,0.0855)(85,0.0817)(120,0.0761)};
        \addlegendentry{$\alpha=0.01$}
        \addplot[MLBlue,thick] coordinates {
          (1,0.3042)(2,0.0995)(3,0.0920)(4,0.0898)(6,0.0860)(8,0.0825)
          (11,0.0775)(15,0.0715)(20,0.0650)(27,0.0574)(36,0.0497)
          (48,0.0422)(64,0.0356)(85,0.0306)(120,0.0270)};
        \addlegendentry{$\alpha=0.1$}
        \addplot[MLAmber,thick] coordinates {
          (1,5.4549)(2,3.8457)(3,2.7171)(4,1.9253)(6,0.9794)(8,0.5124)
          (11,0.2129)(15,0.0886)(20,0.0504)(27,0.0375)(36,0.0316)
          (48,0.0279)(64,0.0260)(85,0.0252)(120,0.0250)};
        \addlegendentry{$\alpha=0.22$}
        \addplot[MLGray,dashed,domain=0:120,forget plot]{0.025};
      \end{axis}
    \end{tikzpicture}
    {\scriptsize Three step sizes that all converge.\par}
  \end{minipage}\hfill
  \begin{minipage}[t]{0.48\linewidth}
    \centering
    \begin{tikzpicture}
      \begin{axis}[mlaxis, width=\linewidth, height=4.8cm, ymode=log,
        xlabel={iteration}, ylabel={$J$},
        xmin=0, xmax=20, ymin=1, ymax=1000,
        legend style={font=\tiny, at={(0.02,0.98)}, anchor=north west}]
        \addplot[MLRose,thick,mark=*,mark size=1pt] coordinates {
          (0,7.75)(2,10.79)(4,15.06)(6,21.03)(8,29.38)(10,41.08)
          (12,57.44)(14,80.34)(16,112.4)(18,157.2)(20,219.9)};
        \addlegendentry{$\alpha=0.25$}
      \end{axis}
    \end{tikzpicture}
    {\scriptsize One that does not: $\alpha>\alpha^\star$.\par}
  \end{minipage}
  \caption{Cost against iteration on the four-point set, from $(0,0)$. The
  dashed line on the left is the minimum $J=0.025$. Note that $\alpha=0.22$
  starts \emph{worse} than $\alpha=0.1$ --- its first step overshoots badly ---
  and yet reaches the minimum first. Bigger is not always slower to settle.}
  \label{fig:alpha-sweep}
\end{figure}

That last observation is worth pausing on, because it contradicts the intuition
students usually arrive with. At $\alpha = 0.22$ the cost after one iteration is
$5.45$, eighteen times worse than the $0.30$ achieved by $\alpha = 0.1$. Yet by
iteration $120$ the aggressive run has reached $0.0250$ while the cautious one
is still at $0.0270$ and falling slowly. Early behaviour is a poor guide.

\subsection{Exactly how large is too large}

For a quadratic cost we can answer this precisely rather than by trial.

\begin{derivation}{The divergence threshold}
Write $\bm{\theta} = (w,b)$ and let $\bm{\theta}^\star$ be the minimiser. Any
quadratic cost can be written
\[
  J(\bm{\theta}) = J^\star
    + \tfrac12(\bm{\theta}-\bm{\theta}^\star)^{\mathsf{T}}
      H\,(\bm{\theta}-\bm{\theta}^\star),
\]
where $H$ is the (constant) matrix of second derivatives. For the squared-error
cost of a linear model,
\[
  H = \frac1m\begin{pmatrix}\sum_i (\xii{i})^2 & \sum_i \xii{i}\\[2pt]
                            \sum_i \xii{i} & m\end{pmatrix}.
\]
The gradient is $\grad J = H(\bm{\theta}-\bm{\theta}^\star)$, so writing
$\mathbf{e} = \bm{\theta}-\bm{\theta}^\star$ for the error, one step of gradient
descent gives
\[
  \mathbf{e} \;\mapsto\; \mathbf{e} - \alpha H\mathbf{e}
              \;=\; (I - \alpha H)\,\mathbf{e}.
\]
Decompose $\mathbf{e}$ along the eigenvectors of $H$. Along the eigenvector with
eigenvalue $\lambda$, the error is multiplied by $(1-\alpha\lambda)$ every step,
so that component shrinks precisely when $|1-\alpha\lambda| < 1$, that is when
$0 < \alpha < 2/\lambda$. All components shrink together only if this holds for
the largest eigenvalue:
\begin{equation}
  \alpha \;<\; \alpha^\star \;=\; \frac{2}{\lambda_{\max}(H)} .
  \label{eq:alpha-star}
\end{equation}
\end{derivation}

For the four-point set,
\[
  H = \begin{pmatrix}7.5 & 2.5\\ 2.5 & 1\end{pmatrix},
  \qquad
  \lambda = 0.1497 \text{ and } 8.3503,
  \qquad
  \alpha^\star = \frac{2}{8.3503} = 0.2395,
\]
which is exactly what \cref{fig:alpha-sweep} shows: $0.22$ converges and $0.25$
does not. The threshold is a property of the \emph{data}, not of the algorithm.
Rescale the feature and it changes; this is one of several reasons
\cref{ch:ch05} rescales features before doing anything else.

\begin{keyinsight}{Both eigenvalues matter, differently}
$\lambda_{\max}$ limits how large $\alpha$ may be. $\lambda_{\min}$ determines
how slowly the last component of the error decays, at rate
$(1-\alpha\lambda_{\min})$ per step. The ratio $\lambda_{\max}/\lambda_{\min}$
--- here $8.3503/0.1497 \approx 56$ --- is therefore what governs the total
number of iterations, and it is the same quantity that makes the contours
elliptical rather than circular.
\end{keyinsight}

\subsection{Choosing $\alpha$ in practice}

You will rarely compute $\lambda_{\max}$. The practical procedure is to try a
ladder of values differing by a factor of about three,
\[
  \dots,\; 0.001,\; 0.003,\; 0.01,\; 0.03,\; 0.1,\; 0.3,\; \dots
\]
running each for a few dozen iterations and plotting $J$. Discard any value for
which $J$ increases; among the rest, take the largest. \Cref{ch:ch05} makes this
routine, and \cref{ch:ch09} replaces it with methods that adapt $\alpha$ on
their own.

\section{Gradient Descent for Linear Regression}
\label{sec:gd-linreg}

Everything so far applies to any differentiable cost. Now specialise.

\begin{derivation}{The two partial derivatives}
With $\fwb(x) = wx+b$ and
$J(w,b) = \frac{1}{2m}\sum_{i=1}^{m}(\fwb(\xii{i})-\yii{i})^2$, differentiate
term by term using the chain rule. For $w$:
\[
  \begin{aligned}
    \frac{\partial J}{\partial w}
      &= \frac{1}{2m}\sum_{i=1}^{m} 2\bigl(w\xii{i}+b-\yii{i}\bigr)
          \cdot\frac{\partial}{\partial w}\bigl(w\xii{i}+b-\yii{i}\bigr)\\[2pt]
      &= \frac{1}{m}\sum_{i=1}^{m}\bigl(\fwb(\xii{i})-\yii{i}\bigr)\,\xii{i}.
  \end{aligned}
\]
For $b$ the inner derivative is $1$ instead of $\xii{i}$:
\[
  \frac{\partial J}{\partial b}
    = \frac{1}{2m}\sum_{i=1}^{m} 2\bigl(w\xii{i}+b-\yii{i}\bigr)\cdot 1
    = \frac{1}{m}\sum_{i=1}^{m}\bigl(\fwb(\xii{i})-\yii{i}\bigr).
\]
This is where the $\tfrac12$ of \cref{ch:ch02} pays for itself: the $2$ from the
chain rule cancels it exactly, leaving no stray constants.
\end{derivation}

Both gradients are averages of the residual $r^{(i)} = \fwb(\xii{i})-\yii{i}$:
weighted by $\xii{i}$ for $w$, unweighted for $b$. That structure is worth
remembering, because it survives every generalisation in this book. In
\cref{ch:ch04} the weighting becomes $\xjj{i}{j}$ for parameter $w_j$; in
\cref{ch:ch06} the residual is computed from a sigmoid instead of a line and
\emph{the formula does not otherwise change}; in \cref{ch:ch09} the same
expression appears at the output layer of a network.

\begin{algorithmbox}{Batch gradient descent for univariate linear regression}
\textbf{Input:} training set $\{(\xii{i},\yii{i})\}_{i=1}^m$, learning rate
$\alpha$, iteration budget $N$.\\
\textbf{Output:} parameters $w,b$.

\begin{enumerate}\setlength{\itemsep}{1pt}
  \item $w \leftarrow 0$, \quad $b \leftarrow 0$.
  \item \textbf{for} $k = 1,\dots,N$:
  \begin{enumerate}\setlength{\itemsep}{1pt}
    \item $g_w \leftarrow \dfrac{1}{m}\sum_{i=1}^{m}(w\xii{i}+b-\yii{i})\,\xii{i}$
    \item $g_b \leftarrow \dfrac{1}{m}\sum_{i=1}^{m}(w\xii{i}+b-\yii{i})$
    \item $w \leftarrow w - \alpha g_w$, \quad $b \leftarrow b - \alpha g_b$
          \hfill \emph{(simultaneously)}
  \end{enumerate}
  \item \textbf{return} $w, b$.
\end{enumerate}

\textbf{Cost per iteration:} $O(m)$ --- every example is visited once.
\textbf{``Batch''} refers to this: each update uses the whole training set.
\end{algorithmbox}

\section{Running It}
\label{sec:running-gd}

\begin{workedexample}{Two iterations by hand on the four-point set}
Data $(1,1),(2,3),(3,4),(4,6)$, so $m=4$. Start at $w=b=0$ with $\alpha = 0.1$.

\textbf{Iteration 1.} With $w=b=0$ every prediction is $0$, so the residuals are
$\hat y^{(i)} - \yii{i} = (-1,\,-3,\,-4,\,-6)$. Then
\[
  g_w = \tfrac14\bigl[(-1)(1)+(-3)(2)+(-4)(3)+(-6)(4)\bigr]
      = \tfrac{-43}{4} = -10.75,
\]
\[
  g_b = \tfrac14\bigl[-1-3-4-6\bigr] = \tfrac{-14}{4} = -3.5 .
\]
Both gradients are negative, so both parameters increase:
\[
  w := 0 - 0.1(-10.75) = 1.0750,
  \qquad
  b := 0 - 0.1(-3.5) = 0.3500 .
\]
The cost has fallen from $J(0,0) = 7.7500$ to $J(1.075,0.35) = 0.3042$.

\textbf{Iteration 2.} Predictions are now $1.425, 2.500, 3.575, 4.650$, so the
residuals are $(+0.425,\, -0.500,\, -0.425,\, -1.350)$, and
\[
  \begin{aligned}
    g_w &= \tfrac14\bigl[0.425 - 1.000 - 1.275 - 5.400\bigr] = -1.8125,\\[2pt]
    g_b &= \tfrac14\bigl[0.425 - 0.500 - 0.425 - 1.350\bigr] = -0.4625 ,
  \end{aligned}
\]
giving $w := 1.2563$, $b := 0.3963$ and $J = 0.0995$.

\textbf{The run continues:}
\begin{center}
\small
\renewcommand{\arraystretch}{1.1}
\begin{tabular}{@{}rrrrrr@{}}
\toprule
$k$ & $w$ & $b$ & $\partial J/\partial w$ & $\partial J/\partial b$ & $J$\\
\midrule
$0$    & $0.0000$ & $ 0.0000$ &            &           & $7.7500$\\
$1$    & $1.0750$ & $ 0.3500$ & $-10.7500$ & $-3.5000$ & $0.3042$\\
$2$    & $1.2563$ & $ 0.3963$ & $ -1.8125$ & $-0.4625$ & $0.0995$\\
$3$    & $1.2900$ & $ 0.3926$ & $ -0.3375$ & $+0.0369$ & $0.0920$\\
$4$    & $1.2994$ & $ 0.3808$ & $ -0.0936$ & $+0.1176$ & $0.0898$\\
$5$    & $1.3046$ & $ 0.3679$ & $ -0.0528$ & $+0.1292$ & $0.0879$\\
$10$   & $1.3262$ & $ 0.3049$ &            &           & $0.0791$\\
$50$   & $1.4503$ & $-0.0597$ &            &           & $0.0412$\\
$100$  & $1.5296$ & $-0.2929$ &            &           & $0.0286$\\
$500$  & $1.5998$ & $-0.4995$ &            &           & $0.0250$\\
\midrule
\multicolumn{2}{@{}l}{optimum (\cref{ch:ch02})} &&&& \\
       & $1.6000$ & $-0.5000$ &            &           & $0.0250$\\
\bottomrule
\end{tabular}
\end{center}

Two features of this table deserve comment. First, $b$'s gradient changes sign
between iterations $2$ and $3$: $b$ overshot upward on the first step and must
now come back down. Individual parameters may move away from their optima; only
$J$ is guaranteed to fall. Second, look at how the work is distributed.
Ninety-six per cent of the reduction in $J$ happens in the first two iterations;
the remaining five hundred are spent creeping along a nearly flat valley floor
while $b$ travels from $+0.40$ to $-0.50$.
\end{workedexample}

\begin{figure}[htbp]
  \centering
  \begin{tikzpicture}
    \begin{axis}[width=0.66\linewidth, height=5.4cm,
      xlabel={$w$}, ylabel={$b$},
      tick label style={font=\scriptsize}, label style={font=\small},
      view={0}{90}, domain=0.2:3.0, y domain=-3:2,
      colormap={ml}{color=(MLBlue) color=(MLTeal) color=(MLAmber)}]
      \costcontour{}{0.1,0.25,0.6,1.2,2.2,3.6}
        {(30*x^2 + 20*x*y + 4*y^2 - 86*x - 28*y + 62)/8}
      \addplot3[MLRose,thick,mark=*,mark size=1.1pt] coordinates {
        (0.000,0.000,0)(1.075,0.350,0)(1.256,0.396,0)(1.290,0.393,0)
        (1.299,0.381,0)(1.305,0.368,0)(1.309,0.355,0)(1.318,0.330,0)
        (1.326,0.305,0)(1.338,0.269,0)(1.350,0.235,0)(1.365,0.192,0)
        (1.382,0.142,0)(1.398,0.095,0)(1.426,0.012,0)(1.450,-0.060,0)
        (1.471,-0.121,0)};
      \addplot3[only marks,mark=star,mark size=3pt,MLBlue]
        coordinates {(1.6,-0.5,0)};
      \node[font=\tiny,MLGray,anchor=north west] at (axis cs:0.25,0.1,0)
        {start};
      \node[font=\tiny,MLBlue,anchor=west] at (axis cs:1.68,-0.5,0)
        {optimum};
    \end{axis}
  \end{tikzpicture}
  \caption{The same run, drawn on the contour map of \cref{fig:contours}. The
  first step crosses most of the parameter space; every subsequent step is short
  and nearly parallel to the valley floor. Steps are perpendicular to the
  contours they start on, which is why an elongated valley is traversed so
  inefficiently --- the direction of steepest descent is not the direction of
  the minimum.}
  \label{fig:gd-path}
\end{figure}

\Cref{fig:gd-path} is the geometric explanation of the table. The negative
gradient points along the steepest descent from wherever you stand, which on a
long narrow valley means almost straight across it rather than along it. The
algorithm therefore makes rapid progress in one direction, then crawls. This is
not a defect of the implementation; it is what steepest descent does on
elliptical contours, and the only cures are to reshape the ellipses
(\cref{ch:ch05}) or to use a method with a memory of previous steps
(\cref{ch:ch09}).

\section{Convexity: Why This Always Works Here}
\label{sec:convexity}

Gradient descent finds a point where the gradient vanishes. In general that is
no guarantee of anything: a function may have many local minima, and which one
you reach depends on where you started. For the squared-error cost of a linear
model, this worry does not arise.

\begin{derivation}{The cost is convex}
A twice-differentiable function is \term{convex} when its matrix of second
derivatives is positive semi-definite everywhere. Differentiating the gradients
of \cref{sec:gd-linreg} once more:
\[
  H = \begin{pmatrix}
        \partial^2 J/\partial w^2 & \partial^2 J/\partial w\,\partial b\\
        \partial^2 J/\partial b\,\partial w & \partial^2 J/\partial b^2
      \end{pmatrix}
    = \frac1m\begin{pmatrix}\sum_i (\xii{i})^2 & \sum_i \xii{i}\\
                            \sum_i \xii{i} & m\end{pmatrix},
\]
which does not depend on $(w,b)$ at all. For any vector $\mathbf{u} = (u_1,u_2)$,
\[
  \mathbf{u}^{\mathsf{T}} H \mathbf{u}
    = \frac1m\sum_{i=1}^{m}\bigl(u_1\xii{i} + u_2\bigr)^2 \;\ge\; 0,
\]
because it is a sum of squares. So $H$ is positive semi-definite and $J$ is
convex; it is strictly convex, with a unique minimum, unless every $\xii{i}$ is
identical. For the four-point set, $H = \left(\begin{smallmatrix}7.5 & 2.5\\ 2.5
& 1\end{smallmatrix}\right)$, with trace $8.5>0$, determinant $1.25>0$ and
eigenvalues $0.1497$ and $8.3503$, both positive.
\end{derivation}

\begin{keyinsight}{What convexity guarantees, and what it does not}
For a convex cost with $\alpha < \alpha^\star$, gradient descent converges to the
global minimum from any starting point. There are no local minima to be trapped
in, no saddle points, no dependence on initialization, and no reason to run it
several times from different starts.

None of that survives \cref{ch:ch09}. A neural network's cost is not convex, its
minima are many, and where you start matters. What does survive is the update
rule, unchanged --- which is why it is worth understanding thoroughly in the
setting where it provably works.
\end{keyinsight}

\section{When to Stop}
\label{sec:stopping}

The definition says ``repeat until convergence'', which is not something a
program can execute. Three rules are in common use:

\begin{enumerate}
  \item \textbf{A fixed iteration budget.} Simple, predictable, and used more
        often than textbooks admit.
  \item \textbf{A small change in cost:} stop when
        $|J_{k-1} - J_k| < \varepsilon$.
  \item \textbf{A small gradient:} stop when $\|\grad J\| < \varepsilon$. This
        is the condition that actually characterises a minimum, and it is the
        most defensible of the three.
\end{enumerate}

The second is the most popular and the least trustworthy, because a small change
in $J$ means the algorithm is moving slowly, which is not the same as having
arrived.

\begin{pitfall}{A flat curve is not a converged run}
Apply rule 2 to the run in \cref{sec:running-gd}, at four thresholds:

\begin{center}
\small
\begin{tabular}{@{}lrrr@{}}
\toprule
$\varepsilon$ & stops at iteration & reported $J$ & true minimum\\
\midrule
$10^{-2}$ & $3$   & $0.0920$ & $0.0250$\\
$10^{-3}$ & $27$  & $0.0574$ & $0.0250$\\
$10^{-4}$ & $104$ & $0.0282$ & $0.0250$\\
$10^{-6}$ & $256$ & $0.0250$ & $0.0250$\\
\bottomrule
\end{tabular}
\end{center}

At $\varepsilon = 10^{-2}$ the rule halts after three iterations, reporting a
cost nearly four times the minimum, and it does so because the run genuinely has
slowed --- the valley floor of \cref{fig:gd-path} is almost flat in the
direction of travel. Nothing about the curve looks wrong. Always plot $J$
against iteration on a \emph{logarithmic} vertical axis, where the difference
between $0.09$ and $0.025$ is plainly visible, and prefer a gradient-norm test
when you can afford it.
\end{pitfall}

\section{Batch, Stochastic and Mini-Batch}
\label{sec:sgd}

Each iteration of the algorithm as stated costs $O(m)$, because both gradients
are averages over the whole training set. When $m$ is ten million, a single
parameter update requires ten million multiply-adds --- and hundreds of updates
are needed. The three standard variants trade accuracy of the gradient against
its cost.

\begin{center}
\small
\renewcommand{\arraystretch}{1.25}
\begin{tabular}{@{}L{2.6cm}L{2.9cm}L{5.3cm}@{}}
\toprule
\textbf{Variant} & \textbf{Examples per update} & \textbf{Character}\\
\midrule
Batch & all $m$ &
  Exact gradient; smooth monotone decrease in $J$; one update per pass over the
  data.\\
Stochastic (SGD) & $1$ &
  Very noisy gradient; $J$ fluctuates and never settles; $m$ updates per pass,
  so early progress is dramatic.\\
Mini-batch & $32$ to $512$, typically &
  A compromise, and what is almost always used in practice; vectorizes well on
  the hardware of \cref{ch:ch04}.\\
\bottomrule
\end{tabular}
\end{center}

The key observation, due originally to \textcite{robbins1951} in a different
context, is that the gradient computed from a single random example is
\emph{unbiased}: its expectation over the choice of example is the true
gradient. The steps are wrong individually and right on average, and averaging
happens automatically as the iterations accumulate.

\begin{figure}[htbp]
  \centering
  \begin{tikzpicture}
    \begin{axis}[width=0.66\linewidth, height=5.2cm,
      xlabel={$w$}, ylabel={$b$},
      tick label style={font=\scriptsize}, label style={font=\small},
      view={0}{90}, domain=0.2:3.0, y domain=-3:2,
      colormap={ml}{color=(MLBlue) color=(MLTeal) color=(MLAmber)}]
      \costcontour{}{0.1,0.25,0.6,1.2,2.2,3.6}
        {(30*x^2 + 20*x*y + 4*y^2 - 86*x - 28*y + 62)/8}
      \addplot3[MLRose,thick,mark=*,mark size=1pt] coordinates {
        (0.000,0.000,0)(0.620,0.240,0)(1.291,0.446,0)(1.349,0.455,0)
        (1.238,0.395,0)(1.369,0.404,0)(1.286,0.374,0)(1.254,0.364,0)
        (1.339,0.356,0)(1.397,0.345,0)(1.321,0.320,0)(1.319,0.313,0)
        (1.302,0.289,0)(1.367,0.278,0)(1.322,0.265,0)(1.311,0.245,0)
        (1.410,0.267,0)};
      \addplot3[only marks,mark=star,mark size=3pt,MLBlue]
        coordinates {(1.6,-0.5,0)};
    \end{axis}
  \end{tikzpicture}
  \caption{Stochastic gradient descent on the four-point set, $\alpha=0.05$, one
  example per step. Compare \cref{fig:gd-path}: the same general journey, taken
  by a drunk. The noise is not a defect to be eliminated --- on a non-convex
  cost it helps the search escape poor regions --- but it does mean SGD never
  settles at the minimum, only wanders near it.}
  \label{fig:sgd-path}
\end{figure}

\begin{researchspotlight}{Robbins \& Monro (1951) --- A Stochastic Approximation Method}
\textbf{The problem.} Find the root of a function $M(\theta)$ --- equivalently,
the minimum of its integral --- when $M$ cannot be evaluated exactly, only
observed with random noise. This arises constantly in experimental work: each
measurement of a response is contaminated, and measurements are expensive.

\textbf{The idea.} Iterate $\theta_{k+1} = \theta_k - a_k Y_k$, where $Y_k$ is a
single noisy observation at $\theta_k$ and $a_k$ is a step size that shrinks over
time. Do not average many observations at each point; take one, step, and let
the iteration do the averaging.

\textbf{The result.} The sequence converges in probability to the true root
provided the step sizes satisfy $\sum_k a_k = \infty$ and $\sum_k a_k^2 <
\infty$: large enough in total to reach any starting distance, small enough in
square to damp the noise. The canonical choice is $a_k \propto 1/k$.

\textbf{Why it is here.} Replace ``noisy observation of the gradient'' with
``gradient computed from one training example'' and this is stochastic gradient
descent, forty years early. The conditions on $a_k$ are the theoretical basis of
learning-rate schedules, and the reason a constant $\alpha$ leaves SGD wandering
rather than converging, exactly as in \cref{fig:sgd-path}.
\parencite{robbins1951}
\end{researchspotlight}

\begin{labbox}{Lab 3 --- Gradient descent from scratch}
Implement \code{gradient\_descent(x, y, alpha, iters)} returning the history of
$(w, b, J)$, using a simultaneous update. Reproduce the iteration table of
\cref{sec:running-gd} exactly --- your iteration 1 must give
$(1.0750,\,0.3500)$ with $J = 0.3042$. Then plot $J$ against iteration for
$\alpha \in \{0.01, 0.1, 0.22, 0.25\}$ and confirm that the last diverges, and
that the threshold $\alpha^\star = 0.2395$ predicted by
\cref{eq:alpha-star} is where the behaviour changes. Finally, deliberately
introduce the non-simultaneous bug and check that you reproduce the wrong
column of the pitfall table --- knowing what the bug looks like is the only
reliable way to notice it later.
\end{labbox}

\section*{Chapter Summary}
\addcontentsline{toc}{section}{Chapter Summary}

\begin{itemize}
  \item Gradient descent minimises a cost using only its value and derivatives,
        which is all that is available for most models. It trades exactness for
        generality.
  \item The update is $w := w - \alpha\,\partial J/\partial w$ for each
        parameter, with all parameters updated \emph{simultaneously}. The
        derivative supplies both direction (its sign) and step length (its
        magnitude), so steps shorten automatically near a minimum.
  \item A sequential update is a different algorithm. It may still reduce $J$,
        so monitoring the cost will not reveal the mistake.
  \item Too small an $\alpha$ is slow; too large diverges. For a quadratic cost
        the exact threshold is $\alpha^\star = 2/\lambda_{\max}(H)$, a property
        of the data. The ratio $\lambda_{\max}/\lambda_{\min}$ governs how many
        iterations are needed.
  \item For squared error, $\partial J/\partial w = \frac1m\sum_i r^{(i)}\xii{i}$
        and $\partial J/\partial b = \frac1m\sum_i r^{(i)}$: averages of the
        residual, weighted and unweighted. The $\tfrac12$ in $J$ cancels the $2$
        from the chain rule.
  \item The cost of a linear model is convex, because its Hessian is a sum of
        squares. Gradient descent therefore reaches the global minimum from any
        start. Neither statement survives into neural networks; the update rule
        does.
  \item Steps are perpendicular to the contours, so elongated valleys are
        crossed quickly and traversed slowly. Most of the reduction in $J$
        happens in the first few iterations; most of the iterations happen
        afterwards.
  \item Stopping when $J$ stops changing much can halt far from the minimum.
        Prefer a gradient-norm test, and always plot $J$ on a logarithmic axis.
  \item Batch, stochastic and mini-batch differ only in how many examples each
        update uses. The single-example gradient is noisy but unbiased;
        mini-batch is the practical default.
\end{itemize}

\section*{Exercises}
\addcontentsline{toc}{section}{Exercises}

\exercisesection{Short Questions}

\sq Write the gradient descent update rule for a cost $J(w,b)$, and say what
$:=$ means.

\sq Why is the derivative \emph{subtracted} rather than added?

\sq Explain, without using the word ``gradient'', why gradient descent takes
shorter steps as it approaches a minimum.

\sq What is a simultaneous update, and why does it matter?

\sq Give one symptom of a learning rate that is too small and two symptoms of
one that is too large.

\sq State the divergence threshold for gradient descent on a quadratic cost, and
say which property of the data determines it.

\sq Write $\partial J/\partial w$ and $\partial J/\partial b$ for the
squared-error cost of a linear model, and identify the common structure.

\sq What does it mean for a function to be convex, and what does convexity
guarantee about gradient descent?

\sq Give two stopping rules, and say which is more trustworthy and why.

\sq Distinguish batch, stochastic and mini-batch gradient descent by the number
of examples per update, and give the cost per update of each.

\exercisesection{Analytical Problems}

\ap For the three-point set $(1,1),(2,2),(3,3)$ with $b$ fixed at $0$, start at
$w = 3$ with $\alpha = 0.15$ and carry out three iterations by hand, reporting
$w$ and $J(w)$ at each step. Does the sequence overshoot the minimum? Compute
$\alpha^\star$ for this problem and check your answer is consistent with it.

\ap Carry out iteration $3$ of the worked example in \cref{sec:running-gd} by
hand, starting from $(1.2563, 0.3963)$, and confirm the printed values
$g_w = -0.3375$, $g_b = +0.0369$ and $J = 0.0920$. Explain in one sentence why
$g_b$ has changed sign since iteration $1$.

\ap Show that for the squared-error cost of a linear model, the gradient can be
written $\grad J = H(\bm{\theta} - \bm{\theta}^\star)$ with $H$ the Hessian and
$\bm{\theta}^\star$ the minimiser. Hence show that one step of gradient descent
multiplies the error by $(I - \alpha H)$.

\ap A dataset has $\sum_i (\xii{i})^2 = 500$, $\sum_i \xii{i} = 60$ and $m = 10$.
Compute $H$, its eigenvalues, and $\alpha^\star$. If the feature is then divided
by $10$, recompute all three. What does this tell you about choosing $\alpha$
before scaling features?

\ap Suppose $J(w) = \tfrac12 \lambda w^2$ for some $\lambda > 0$. Show that
gradient descent from $w_0$ gives $w_k = (1-\alpha\lambda)^k w_0$, and hence
determine the value of $\alpha$ that converges fastest. What happens at exactly
that value, and at twice it?

\ap Prove that if $J$ is convex and $\alpha$ small enough that each step
decreases $J$, then the sequence $J_0, J_1, J_2, \dots$ is monotone decreasing
and bounded below, and hence converges. Does it follow that the parameters
converge?

\exercisesection{Design Problems}

\dpr You are given a training run in which $J$ decreases for $200$ iterations and
then increases sharply. List the possible causes in order of likelihood, and for
each say what you would change and what you would expect to see if that cause
was the right one.

\dpr A colleague's model trains for six hours. They propose increasing $\alpha$
tenfold to save time. Set out what you would check before agreeing, what
evidence would persuade you, and what cheaper interventions you would try
first --- referring to \cref{ch:ch05} where appropriate.

\dpr Design a stopping rule for a training job that must finish within a fixed
wall-clock budget while producing the best model it can. State what it monitors,
what it stores, and how it behaves if the budget expires mid-run.

\dpr Your team trains on ten million examples. Argue for batch, stochastic, or
mini-batch gradient descent, naming the batch size you would choose and the
evidence you would collect to justify it. What would change your mind?

\exercisesection{Programming Exercises}

\pe Implement batch gradient descent for the univariate model and reproduce the
iteration table of \cref{sec:running-gd} to four decimal places.

\pe Plot $J$ against iteration on a logarithmic axis, on the four-point set, for
each of the step sizes $0.01$, $0.05$, $0.1$, $0.22$, $0.24$ and $0.25$. Find
empirically the largest $\alpha$ that still converges, to two decimal places,
and compare it with $\alpha^\star = 0.2395$.

\pe Record the parameter path and overlay it on a contour plot of $J$, as in
\cref{fig:gd-path}. Then rescale the feature by subtracting its mean and
dividing by its standard deviation, refit, and overlay the new path on the new
contours. Comment on the number of iterations required in each case. (This is
\cref{ch:ch05} in advance.)

\pe Implement stochastic gradient descent and plot $J$ against the number of
examples processed, for both batch and stochastic versions on the same axes.
Which reaches $J < 0.03$ after fewer example-visits? Then implement a decaying
learning rate $\alpha_k = \alpha_0/(1+k/100)$ and show that SGD now settles
rather than wandering.

\exercisesection{Self-Test}

\qz In the update $w := w - \alpha\,\partial J/\partial w$, the role of the
derivative's sign is to determine: (a)~the step length; (b)~the direction;
(c)~the stopping point; (d)~the convexity.

\qz Gradient descent with too large a learning rate typically shows:
(a)~$J$ decreasing very slowly; (b)~$J$ increasing or oscillating;
(c)~$J$ constant; (d)~a negative $J$.

\qz The divergence threshold $2/\lambda_{\max}$ depends on: (a)~the starting
point; (b)~the training data; (c)~the number of iterations; (d)~the initial
cost.

\qz The gradient $\partial J/\partial b$ for squared error equals: (a)~the sum
of the residuals; (b)~the mean of the residuals; (c)~the mean residual weighted
by $x$; (d)~zero.

\qz Updating $w$ and then computing $b$'s gradient with the new $w$: (a)~is
equivalent to the simultaneous update; (b)~always diverges; (c)~is a different
algorithm that may still reduce $J$; (d)~halves the learning rate.

\qz The cost of a linear model with squared error is convex because its Hessian:
(a)~is diagonal; (b)~is a sum of squares and so positive semi-definite;
(c)~depends on $(w,b)$; (d)~has determinant zero.

\qz Gradient descent steps are: (a)~parallel to the contours; (b)~perpendicular
to the contours; (c)~directed at the minimum; (d)~random.

\qz Compared with batch gradient descent, one step of stochastic gradient
descent is: (a)~more accurate and more expensive; (b)~less accurate and cheaper;
(c)~equally accurate and cheaper; (d)~biased and cheaper.

\furtherreading

\textcite{james2021} treat linear regression through the closed form and say
little about optimization, so for this chapter the primary reference is
§3.1 for context only. \textcite{bishop2006} §3.1.3 introduces sequential
(stochastic) gradient descent for exactly this cost function and connects it to
the Robbins--Monro conditions, and §5.2.4 discusses learning rates and the
geometry of ill-conditioned error surfaces --- the analytical version of
\cref{fig:gd-path}.

\bigskip
\noindent\emph{Next:} \Cref{ch:ch04} lets the model see more than one feature.
The update rule survives untouched; the notation, the arithmetic cost, and the
opportunity for a hundredfold speed-up all change.
